
\documentclass[12pt]{article}
\usepackage[margin=1in]{geometry}
\usepackage{setspace}
\usepackage{amsmath, amssymb}
\usepackage{natbib}
\usepackage{hyperref}

\title{Political Polarization and the Interpretation of Auditor Signals}
\author{Neil}
\date{\today}

\begin{document}
\maketitle
\onehalfspacing

\section{Introduction}

Do investors interpret auditor signals uniformly? The auditing literature largely assumes that audit reports and related disclosures convey information that is processed similarly across investors. In contrast, a growing literature in finance and economics shows that political polarization shapes beliefs, information processing, and trading behavior. This paper bridges these literatures by examining whether political polarization affects how investors interpret auditor-related information.

We argue that political polarization induces heterogeneity in prior beliefs and trust in institutions, leading to differential interpretation of identical audit signals. As a result, polarization affects not only price reactions but also disagreement, trading volume, and volatility around audit-related events.

\section{Related Literature}

\subsection{Audit and Capital Markets}

Prior work establishes that audit reports contain information relevant to investors \citep{firth1978,dodd1984}. Subsequent literature studies audit quality and market outcomes \citep{francis2004,defond2014}. Research on expanded audit reporting finds mixed evidence on incremental informativeness.

\subsection{Political Economy of Auditing}

A related literature examines how political connections affect audit outcomes \citep{chan2006}. However, these studies focus on auditor behavior rather than investor interpretation.

\subsection{Political Polarization in Finance}

Recent research shows that political polarization affects financial decisions and asset pricing \citep{kempf2024}. Investors exhibit partisan behavior, and political beliefs influence trading and expectations.

\subsection{Disagreement and Information Processing}

Theoretical models show that heterogeneous beliefs lead to trading and volatility \citep{kandel1995,harris1993,hong2007}. We extend this framework by linking belief heterogeneity to political polarization.

\section{Theory}

We consider a setting where investors observe a public audit signal $s$ about firm quality $\theta$. Investors differ in priors $\pi_i$ due to political polarization.

\begin{equation}
P_i(\theta=1|s) = \frac{P(s|\theta=1)\pi_i}{P(s)}
\end{equation}

Polarization increases dispersion in $\pi_i$, leading to greater dispersion in posterior beliefs.

\subsection{Equilibrium}

Following \citet{harris1993}, prices reflect weighted average beliefs:

\begin{equation}
P = \sum_i w_i E_i[\theta|s]
\end{equation}

Greater dispersion in beliefs increases trading volume and absolute returns.

\subsection{Testable Implications}

\begin{itemize}
\item Polarization increases absolute returns around audit events.
\item Polarization increases trading volume.
\item Effects are stronger for ambiguous audit signals.
\end{itemize}

\section{Empirical Design}

\subsection{Setting}

We study auditor changes disclosed in Form 8-K (Item 4.01).

\subsection{Specification}

\begin{equation}
CAR_{i,t} = \alpha + \beta_1 Event_{i,t} + \beta_2 Polarization_{i,t} + \beta_3 Event \times Polarization + FE + \varepsilon
\end{equation}

\subsection{Variables}

\textbf{Dependent variables:} CAR, absolute CAR, trading volume.

\textbf{Key independent variables:} auditor change indicators, polarization measures.

\section{Identification}

We exploit cross-sectional variation in polarization and event-level shocks. Firm and time fixed effects mitigate omitted variable concerns.

\section{Conclusion}

This paper shows that political polarization affects how investors interpret audit signals, contributing to auditing, finance, and political economy literatures.

\bibliographystyle{apalike}
\bibliography{references}

\end{document}
